% Main blueprint content. Each \begin{theorem}/\begin{definition} block names
% one or more Lean declarations via \lean{...} and may declare dependencies
% via \uses{...}. Mark fully-formalized statements/proofs with \leanok.

\chapter{The Four Color Theorem}

\section{Top-level statement}

\begin{definition}[Simple map]
  \label{def:simple-map}
  \lean{FourColor.RealPlane.SimpleMap}
  \leanok
  A \emph{simple map} of the real plane is a partial equivalence relation on
  $\R^2$ whose classes (the regions) are open and connected.
\end{definition}

\begin{definition}[$n$-colorable]
  \label{def:colorable-with}
  \lean{FourColor.RealPlane.ColorableWith}
  \leanok
  \uses{def:simple-map}
  A map $m$ is $n$-\emph{colorable} if there is a coloring of its regions by
  one of $n$ colors such that adjacent regions receive different colors.
\end{definition}

\begin{theorem}[Four Color Theorem]
  \label{thm:four-color}
  \lean{FourColor.four_color}
  \uses{def:simple-map, def:colorable-with, thm:four-color-finite,
        thm:compactness-extension}
  Every simple map of the real plane is $4$-colorable.
\end{theorem}

\begin{proof}
  \uses{thm:four-color-finite, thm:compactness-extension}
  Combine the finite case (\cref{thm:four-color-finite}) with the compactness
  extension (\cref{thm:compactness-extension}).
\end{proof}

\begin{theorem}[Four Color Theorem, finite case]
  \label{thm:four-color-finite}
  \lean{FourColor.four_color_finite}
  \uses{thm:four-color-hypermap, thm:discretize-to-hypermap}
  Every finite simple map is $4$-colorable.
\end{theorem}

\begin{proof}
  \uses{thm:four-color-hypermap, thm:discretize-to-hypermap}
  Discretize the map to a planar bridgeless hypermap
  (\cref{thm:discretize-to-hypermap}) and apply the combinatorial 4CT
  (\cref{thm:four-color-hypermap}).
\end{proof}

\section{Combinatorial core}

\begin{theorem}[Combinatorial Four Color Theorem]
  \label{thm:four-color-hypermap}
  \lean{FourColor.four_color_hypermap}
  \uses{thm:planar-cube, thm:cube-colorable, thm:unavoidability,
        thm:reducibility}
  Every planar bridgeless hypermap is four-colorable.
\end{theorem}

\section{Walkup construction and Jordan equivalence}

\begin{theorem}[WalkupE Euler components]
  \label{thm:walkupE-euler-components}
  \lean{FourColor.walkupE_euler_components}
  Records how the Euler-formula components $\bigl(\#E,\#N,\#F,\#\mathrm{Comp}\bigr)$
  change when a dart is removed by the WalkupE construction.
\end{theorem}

\begin{theorem}[Jordan WalkupE]
  \label{thm:jordan-walkupE}
  \lean{FourColor.jordan_walkupE}
  \uses{thm:walkupE-euler-components}
  The Jordan property (no Möbius paths) is preserved by WalkupE.
\end{theorem}

\begin{theorem}[Jordan implies planar]
  \label{thm:jordan-planar}
  \lean{FourColor.jordan_planar}
  \uses{thm:jordan-walkupE}
\end{theorem}

\begin{theorem}[Planar implies Jordan]
  \label{thm:planar-jordan}
  \lean{FourColor.planar_jordan}
  \uses{thm:jordan-walkupE}
\end{theorem}

\section{Cube construction}

\begin{theorem}[Face orbit count of the cube]
  \label{thm:fcardFace-cube}
  \lean{FourColor.fcardFace_cube}
  $\#F(\mathrm{cube}\,G) = \mathrm{eulerRhs}(G)$.
\end{theorem}

\begin{theorem}[Cube preserves planarity]
  \label{thm:planar-cube}
  \lean{FourColor.planar_cube}
  \leanok
  \uses{thm:fcardFace-cube}
\end{theorem}

\begin{theorem}[Cube preserves four-colorability]
  \label{thm:cube-colorable}
  \lean{FourColor.cube_colorable}
  \leanok
\end{theorem}

\section{Configurations, reducibility and unavoidability}

\begin{definition}[The 633 configurations]
  \label{def:theConfigs}
  \lean{FourColor.theConfigs}
  The list of 633 reducible configurations (Gonthier et al.).
\end{definition}

\begin{theorem}[Configurations count]
  \label{thm:theConfigs-length}
  \lean{FourColor.theConfigs_length}
  \uses{def:theConfigs}
  $|\mathrm{theConfigs}| = 633$.
\end{theorem}

\begin{definition}[C-reducibility]
  \label{def:CfReducible}
  \lean{FourColor.CfReducible}
\end{definition}

\begin{theorem}[Reducibility]
  \label{thm:reducibility}
  \lean{FourColor.the_reducibility}
  \uses{def:theConfigs, thm:theConfigs-length, def:CfReducible}
  All 633 configurations are C-reducible.
\end{theorem}

\begin{theorem}[Unavoidability]
  \label{thm:unavoidability}
  \lean{FourColor.unavoidability}
  \uses{thm:reducibility}
  Reducibility implies that no minimal counter-example exists.
\end{theorem}

\section{Discretization and compactness}

\begin{theorem}[Discretize to hypermap]
  \label{thm:discretize-to-hypermap}
  \lean{FourColor.discretize_to_hypermap}
  Every finite simple map admits a planar bridgeless hypermap whose
  four-colorings induce four-colorings of the original map.
\end{theorem}

\begin{theorem}[Compactness extension]
  \label{thm:compactness-extension}
  \lean{FourColor.compactness_extension}
  If every finite simple map is $n$-colorable, then every simple map is
  $n$-colorable.
\end{theorem}
